% ============================================================
%  Poul Martin Møller,
%  "Ontologien eller Kategoriernes System"
%
%  A fragment (Brudstykke) on ontology / the system of the
%  categories. Per the author's footnote, "Ontologien eller
%  Kategoriernes System" is the name Møller gave the science he
%  lectured on in the winter half-year of 1837, announced in the
%  review (Anmeldelse) of those lectures.
%
%  Published posthumously in:
%  Efterladte Skrifter af Poul M. Møller, bd. 3 (Tredie Bind),
%  ed. F.C. Olsen. Kjøbenhavn: C.A. Reitzel.
%  (Transcribed from the 3rd edition, "Tredie Udgave".)
%  Printed pp. 189--212. PDF (bibliotek scan) pp. 200--223;
%  offset PDF = printed + 11.
%
%  TRANSCRIPTION (Danish)
%  Conventions: original orthography kept (aa-spellings,
%  capitalised nouns, æ/ø, "philosophisk", "Categorie/Kategorie"
%  as printed). Fraktur long-s rendered as s. Letterspaced
%  emphasis in the original -> \emph{}. Printed page numbers ->
%  \opage{N} (marginal, at the point the printed page begins).
% ============================================================

\documentclass[12pt]{article}
\usepackage[utf8]{inputenc}
\usepackage[T1]{fontenc}
\usepackage{graphicx}
% Old Danish "id est" mark (reversed c + colon), printed as "ɔ:"
\DeclareUnicodeCharacter{0254}{\reflectbox{c}}
\usepackage{libertinus}
\usepackage{libertinust1math}
\usepackage{textalpha}
\usepackage[danish]{babel}
\usepackage[protrusion=true,expansion=true]{microtype}
\usepackage[margin=1.4in]{geometry}
\usepackage{setspace}
\usepackage{fancyhdr}
\usepackage{marginnote}
\usepackage[colorlinks=true,linkcolor=black,urlcolor=black,
  pdftitle={Ontologien eller Kategoriernes System},
  pdfauthor={Poul Martin Møller}]{hyperref}
\setlength{\headheight}{15pt}
\pagestyle{fancy}
\fancyhf{}
\fancyhead[LE,RO]{\thepage}
\fancyhead[RE]{\textit{Ontologien eller Kategoriernes System}}
\fancyhead[LO]{\textit{Poul Martin Møller}}
\setlength{\parindent}{1.5em}
\setlength{\parskip}{0pt}
\onehalfspacing

\newcommand{\opage}[1]{\marginnote{\footnotesize\textit{[#1]}}}
\newcommand{\ombreak}{\par\medskip\begin{center}\rule{0.32\textwidth}{0.4pt}\end{center}\medskip}

\title{\textbf{Ontologien}\\[0.3em]
  \large eller\\[0.3em]
  \textbf{Kategoriernes System}\footnote{Saaledes betegner Forf.\ den
  foredragne Videnskab i Anmeldelsen af hans Forelæsninger over den i
  Vinterhalvaaret 1837.}}
\author{Poul Martin Møller}
\date{Efterladte Skrifter, bd.\ 3}

\begin{document}
\maketitle

\begin{center}\textsc{Indledning.}\end{center}

\opage{189}Indledningen til Metaphysiken i Ordets egentlige Betydning er en
særegen philosophisk Disciplin, der udvikler de forskjellige Stadier, som den
menneskelige Bevidsthed gjennemløber fra sin umiddelbare Skikkelse til det
Dannelsestrin, hvor Metaphysiken tager sin Begyndelse. En saadan philosophisk
Propædeutik har Hegel givet i sin Phänomenologie des Geistes, men dette Værk
overskrider ved sin Mangfoldighed af Indhold langt de ved hiin Definition
betegnede Grændser og forudsætter desuden Læsere, der have Bekjendtskab med
Philosophiens Historie. Ogsaa andre Philosopher have under forskjellige
Benævnelser givet Fremstillinger af den philosophiske Propædeutik i hiin
Betydning. De meest udmærkede skyldes J.\ E.\ von Berger, C.\ F.\ Krause og
J.\ H.\ Fichte. I en vis Betydning kan da Metaphysiken ei kaldes Philosophiens
Begyndelse, da Propædeutiken i et hensigtsmæssigt Foredrag af Philosophiens
System maa forudskikkes. Denne Viden\opage{190}skab skal nemlig lede Læseren
til Metaphysikens Grændser og saaledes orientere ham, at han klart erkjender, i
hvilken Region han nu skal føres ind. Den skal lære ham at forstaae, hvad der
menes, naar man kalder den Tænkning, han nu bliver opfordret til at følge,
immanent eller apriorisk Tænkning. Med Hensyn til Læreforedragets naturlige
Orden er da Propædeutiken at betragte som den første Disciplin af Philosophiens
System; efter et andet Hensyn maa Metaphysiken kaldes Philosophiens Begyndelse.
Den gaaer nemlig ud fra de almindeligste Tankebestemmelser for al Virkelighed,
der maae forudsættes i de concretere Begreber, der gjennemgaaes i
Propædeutiken. I Begrebet er altsaa Metaphysiken den sandere Begyndelse, hvortil
hiin forberedende Videnskab gaaer tilbage.

Metaphysiken kan da ei alene, som enhver philosophisk Disciplin, paa Grund af
sin relative Selvstændighed foredrages saaledes, at den for sig taget bliver
forstaaelig, men den kan med større Ret, end nogen af Philosophiens øvrige Dele,
i sin isolerede Skikkelse gjøre Fordring paa Forstaaelighed. Dens Bestemmelser
gjentages nemlig i concretere Skikkelser i alle Philosophiens Dele, men den selv
indeholder de meest enkelte Elementer af al Erkjendelse. Naar Propædeutiken ved
at føre til den Region, hvori den metaphysiske Tænkning bevæger sig, har udført
sin Tjeneste, kan Propædeutiken tabes af Sigte som en Stige, der har opfyldt sin
Bestemmelse. Den nye Tænkning, som man nu begiver sig til, er nemlig ikke længer
betinget ved hiin Indledning, men træder frem som aldeles forudsætningsløs og
umiddelbar. Som en Følge heraf vil den tænkende Læser, der følger Metaphysikens
Foredrag med sin Opmærksomhed (om han end ei strax ved dets Begyndelse mærker, i
hvad Sphære han befinder sig), under dets Fremgang komme til at forstaae det
Lys, som de forskjellige Sætninger ved deres indbyrdes Sammenhæng kaste paa
hinanden.

\opage{191}Paa de faa Sider, som jeg her har givet Navn af Indledning, er det
naturligviis ei min Hensigt at følge den Tankegang, som kræves af en streng
videnskabelig Propædeutik. Jeg forudskikker blot et Par aphoristiske
Bemærkninger, hvorved jeg troer hos en og anden Læser at kunne vække en
foreløbig Forestilling om Metaphysikens Indhold og dens Betydning i Videnskabens
hele System. Om da disse indledende Betragtninger træffe Berøringspuncter i
Læserens Bevidsthedskreds, eller ikke, er aldeles tilfældigt. Det gjelder om
dem, hvad der gjelder om Aphorismer i Almindelighed: de ligne Hatte, der gjøres
uden Maal; det er et Lykketræf, om nogen af dem passer til et vist Individs
Hoved, eller ei.

Hvis man for en Kjender af den kritiske Philosophie skulde fremstille et
foreløbigt Begreb om den Hoveddeel af Metaphysiken, som her skal foredrages,
kunde man give ham den Oplysning derom, at den gik ud paa en systematisk
Deduction af Kategorierne; thi efter et Tillæg af nogle faa nærmere Bestemmelser
vilde hiin Definition nogenlunde give ham en Forestilling om Videnskabens
(Skriftets) Indhold. I den kritiske Philosophie stræbte man nemlig efter at give
en Fortegnelse over de Begreber, som ligge til Grund for Tænkningen, og som ei
ere givne ved den sandselige Fornemmelse. Disse Begreber søgte Kant at bringe
til Bevidsthed ved Reflexion over Dommenes forskjellige Former. For denne
Undersøgelse lagde han den af de sædvanlige Logiker bekjendte Inddeling af
Dommene til Grund (nemlig den, som tager Hensyn til deres Qvalitet, Qvantitet,
Modalitet og Relation). Det blev i hans Philosophie taget som en Kjendsgjerning,
at hine Begreber forudsættes i al Tænkning, der sætter objectiv Realitet. Men
han gjør ikke noget Forsøg paa at vise, hvorledes det ligger i Sagens Natur, at
Tænkningen netop fremtræder i disse Former. Heller ikke gjør han synderlig Rede
for disse Begrebers indbyrdes Forhold til hinanden; da dog naturligen det
Spørgsmaal maatte \opage{192}opkastes, om disse Begreber ere hverandre blot
sideordnede uden videre Sammenhæng, eller om de staae i Forhold til hverandre
som almindelige og særlige. --- Kant giver selv Anledning til dette Spørgsmaal,
idet han i Forbigaaende løselig henkaster den Bemærkning, at den tredie
Kategorie i enhver Klasse fremkommer ved en Forbindelse af den første og den
anden. Men denne løst henkastede Tanke, hvorpaa Kant selv ikke lægger nogen
Vægt, opfordrer stærkt til fremdeles at spørge, hvad Forhold da de forskjellige
Klasser af Kategorier have til hinanden. Slige Gaader udtales meer eller mindre
bestemt alt under den kritiske Philosophies Udvikling, hvor de sammenfattedes i
den bestemte Fordring om en Deduction af Kategorierne.

Men ved hiin Tilflugt til den kritiske Philosophies Talebrug er det af særdeles
Vigtighed, at man slaaer enhver Forestilling om Kategoriernes Subjectivitet af
Hovedet. Kant antog nemlig, at disse Begreber, som ligge til Grund for
Tænkningen, blot ere Former, der ere nødvendige for Mennesket, naar det skal
kunne tillægge sine sandselige Forestillinger objectiv Realitet, uden at disse
derfor som sande Bestemmelser gjelde for Tingen i sig selv. Denne reent
vilkaarlige Paastand erkjendes for ugyldig i en rigtig Fremstilling af
Ontologien.

De Bestemmelser, som under Tænkningen af Tilværelsen medføre fuldkommen
Overbeviisning, gjelde nødvendigviis for Tingen i sig selv. Antager man ikke
denne Sætning for rigtig, indvikles man i en Tvivl, der consequent forfulgt
ophæver sig selv. Man kan vel ikke gjendrive den Skeptiker, der ved ethvert
Skridt af en videnskabelig Tankegang haardnakket holder fast ved den Bemærkning,
at den menneskelige Tænkning er bunden til subjective Love, i Følge hvilke han
danner sig blot skjæve Begreber om Tingenes Natur og Væsen; thi ved ethvert
tænkeligt Beviis behøver han blot at udtale sin eengang for alle færdige Formel:
„Saaledes nødes jeg efter min Subjectivitet til at tænke \opage{193}mig Tingenes
Væsen; men hvo kan indestaae mig for, at jeg ikke, uden at vide det, er bunden
til den Skjebne, at føres til blot vildfarende Tanker, der dog for mig medføre
fuldkommen Overbeviisning.“ --- Videnskaben kan vel udtale en Verdensanskuelse
for hiin absolute Skeptiker, efter hvilken Overbeviisningen om Tænkningens
Realitet styrkes ved et Begreb om Tænkningens Betydning i Tilværelsen og dens
Udspring fra Religionens Gjenstand, det Alt-indbegribende Væsen; men Skeptikeren
kan, naar Begrebet er udtalt, paany fremkomme med sin simple Formular:
„Saaledes maae vi tænke os Tilværelsens Grundforhold, men maaskee de i sig selv
ikke ere saadanne, som vi maae tænke os dem.“ --- Vi nødsages da til at svare
hiin Skeptiker, der troer med een eneste Trylleformel at kunne tilintetgjøre
alle den videnskabelige Tænknings Frugter: Lad det altsaa være muligt, at
Fornuftens Brug fører til blot Blændværk, --- en Mulighed, der for sig tagen
ligesaa lidet kan bestrides, som saa mangen anden, f.\ Ex.\ den, at Maanen er en
gloende Steen --- da philosophere vi under den Forudsætning, at hvad der for en
fornuftig Tænkning medfører Overbeviisning, gjelder for Tingen i sig selv. Vi
antage da eengang for alle, at, hvad fornuftige Væsner nødvendig maae tænke at
være, nødvendig er, saaledes som det maa tænkes; og begjere ingen sikkrere
Videnskab end den, der for os medfører fuldkommen Overbeviisning. --- Efter
disse Indskrænkninger og Berigtigelser kunne vi foreløbig lade hiin Definition
gjelde, at Ontologien er en Deduction af Kategoriernes System.

For en Naturphilosoph af den schellingske Skole forklares Ontologien bedst som
Tydeliggjørelsen (Udviklingen) af det Absolutes Begreb. I Schellings
Philosophie optræder nemlig dette Begreb uden al Indledning, og de faa
Bestemmelser, hvorved det tydeliggjøres, ere reent uforstaaelige for den, der ei
har noget Bekjendtskab med de philosophiske Bevægelser, der gik forud
\opage{194}for Naturphilosophien. Alene ved denne historiske Kundskab kunne
Begreberne: \emph{intellectuel Anskuelse, det Absolute, Subjectivitet,
Objectivitet \&c.}, hvoraf Schelling gjør en umiddelbar Brug, for Læseren have
nogen bestemt Betydning. Denne Mangel har man ogsaa tidt nok betegnet med det
træffende Udtryk, at Metaphysiken eller Ontologien savnedes hos
Naturphilosopherne.

De forskjellige Definitioner paa det Absolute, som findes hist og her i
Schellings mange Skrifter, maatte naturligviis sammenbringes af den opmærksomme
Læser. Det Absolute bliver snart betegnet som Identiteten af det Ideale og
Reale, snart som Identiteten af Eenheden og Modsætningen, snart som Baandet
mellem Eenheden og Mangfoldigheden, snart som Identiteten af Væsen og Form eller
det Bekræftende og det Bekræftede o.\ s.\ v. Men der findes aldeles ingen
tydelig Bevidsthed om disse forskjellige Grundbegrebers Betydning og Forhold til
hinanden. Og dog er det undertiden af Sammenhængen klart, at det er Sagens egen
Natur, der bringer Schelling til i nye Regioner af sit System pludselig at
forandre sin hidtil brugte Terminologie. Til denne Ubestemthed og Ustadighed i
Begrebernes Betegnelse kommer desuden den Mislighed, at Naturphænomeners
Benævnelser jævnlig blive brugte istedenfor rene Tankebestemmelser, som naar
Ordene \emph{Lys} og \emph{Tyngde} ei tages i den sædvanlige Mening, men bruges
for at betegne Fundamentalbegreber i Systemet.

Da Tidsalderen ikke længer fandt sig tilfredsstillet ved den halvpoetiske Form,
hvori Schelling havde fremstillet sin Anskuelse af det Absolutes Natur og Væsen
(især efterat en Sværm af Efterlignere havde bragt Fremstillingens Mangler til
klar Bevidsthed), angreb og udførte G.\ W.\ Hegel det store Værk, at udvikle
alle Philosophiens Grundbegreber i systematisk Sammenhæng. Ved hans Arbeider er
Ontologiens Idee i den Form, som Philosophiens nuværende Skikkelse kræver,
bleven gjenfødt i den philoso\opage{195}phiske Verden. Følgen deraf har ogsaa
været, at de flere udmærkede Philosopher, som i vor Tid have givet
Fremstillinger af denne gjenopvaagnede Videnskab, meer eller mindre have benyttet
hans Forarbeider og dannet sig efter det store Forbilled, som han har opstillet i
sin Wissenschaft der Logik.

Ontologiens Formaal bliver da at bringe alle de Begreber til videnskabelig
Tydelighed, som Philosophien maa betjene sig af i det Absolutes fuldstændige
Definition. Det Absolute skal deri betragtes efter sit Indholds almindelige
Bestemmelser uden Hensyn til de specielle Former, hvorunder det træder frem i
Tid og Rum. Alle de Begreber, som efter deres dybere Betydning kunne tænkes som
væsentlige Prædicater for det Absolute, skulle fattes efter denne Betydning og
afgive Bidrag til en Construction af Erkjendelsen i det Absolutes Lys.

Vi have ovenfor sagt, at Ontologien er en for sig afsluttet Tankekreds, som ved
en klar Fremstilling, ligesom enhver anden philosophisk Disciplin, for sig taget
kan gjøres forstaaelig, uagtet større Klarhed udbredes over den, naar den sees i
sin organiske Sammenhæng med hele Philosophiens System. Men skjøndt Ontologien
har en saadan relativ Selvstændighed, at den kan forstaaes af sig selv, idet
dens forskjellige Dele kaste det nødvendige Lys paa hinanden, er det dog ikke
sagt, at en Læser, der er blottet for al philosophisk Cultur, strax fatter dens
allerførste Sætninger. Hvis en Saadan strax ved Foredragets Begyndelse finder
sig orienteret og mærker, i hvilken Tankens Region han befinder sig, maa det
snarere betragtes som et Lykketræf end som en Virkning, der kan gjøres Regning
paa og som ligger i Sagens Natur.

Men det er ikke sandsynligt, at dette Skrift vil blive læst af ret Mange, der
aldeles mangle Bekjendtskab med Philosophiens Elementer. Tvertimod anseer jeg
det for rimeligt, at de Fleste, som ville skjenke det nogen Opmærksomhed, have
læst eller stu\opage{196}deret Professor Sibberns vigtigste Værker. For
saadanne Læsere forudskikkes endnu nogle faa Betragtninger, hvori blot eet
Begreb lægges til Grund som aldeles bekjendt, jeg mener Begrebet: aprioriske
Begreber eller Kategorier. Det antages her for en afgjort Sag, at der til Grund
for al Erkjendelse ligger endeel Begreber, som ikke fremstille noget ved
sandselig Fornemmelse givet Indhold. Enhver Art af Tænkning i vidtløftig
Betydning af Ordet har ved ethvert Skridt af sin Udvikling saadanne Kategorier
til sine Forudsætninger, og anvender dem, selv uden Bevidsthed derom, ved enhver
af sine Bevægelser. Vi have sagt, at disse Begreber ligge til Grund for al
Erkjendelse. Denne Sætning har gandske bogstavelig Sandhed og fuldkommen
almindelig Gyldighed; thi al Erkjendelse kommer til Virkelighed som Tænkning, og
Tænkningen bestaaer i en Anvendelse af aprioriske Forudsætninger, som udspringe
blot af den selv og lede dens Virksomhed fremad. Ved disse Kategorier erkjendes
Bestemmelser ved Virkeligheden, som kun ere til for Tænkningen, og som dog for
det tænkende Subject medføre samme Overbeviisning om objectiv Gyldighed, som de
Bestemmelser, der i Følge Affectioner af Nervesystemet tillægges Tingene.
Tænkningen gjennemtrænger Erkjendelsen paa ethvert af dens Udviklingstrin: ei
engang en Act af den udvortes Sandsning kan finde Sted uden saaledes, at
Subjectet adskiller noget Tilværende fra sig selv og sætter det som Noget for
sig, forskjelligt fra andet Tilværende; men her komme da strax adskillige
Kategorier til Anvendelse, nemlig Begreberne: Tilværelse, For-sig-Væren o.\ s.\ v.

Der gives i Bevidsthedens Udvikling en Periode, da Mennesket tænker efter sine
Kategorier med fuldkommen Sikkerhed, uden Betænkelighed og Tvivl om deres rette
Betydning. I denne Tænkningens Uskyldighedstilstand udgjøre disse Begreber de
faste Støtter, hvorom Erkjendelsens Udvikling dreier sig, men som selv fritages
for Bevægelsens Uro. Selv efterat Mennesket ved \opage{197}Reflexion over
Tænkningens Forudsætninger og Efterligning af fremmed Tale, der gaaer over dets
Forstand, har tabt hiin sikkre Tact for disses rette Brug, beholde dog en heel
Deel af de meest abstracte Kategorier deres fulde Betydning for det. Det var
heller ikke tænkeligt, at alle Kategorier i et Menneskes Bevidsthed paa eengang
skulde bringes i Forvirring og Opløsning, uden at Forrykthed var den
uudeblivelige Følge deraf.

Alt under hiin Tænkningens blot naturlige Sundhed kan Bevidstheden naturligviis
meer eller mindre svare til sin Idee. Den Tydelighed, hvormed de aprioriske
Begreber paa dette Standpunct komme til Udvikling hos Individet, beroer nærmest
paa den Rigdom og Bestemthed, som hersker i Modersmaalets Ordforraad, og dernæst
paa den Skikkelse, hvorunder Sproget først kommer det for Øren i dets
Omgangskreds. Den, der bliver vakt til Selvbevidsthed ved et Sprog, der er saa
fattigt paa Tankebestemmelser, som det grønlandske, bliver ved en uundgaaelig
Skjebne sat paa et lavt Trin af intellectuel Cultur, som hans individuelle
Naturgaver ikke formaae at hæve ham stort over.

Vi tale her om Kategorierne, saaledes som de først vise sig, naar de blot
umiddelbar fastholdes ved bestemte Former i Sproget. Da ere de nærmest at
betragte som en Tradition, som en Folkestammes fælleds Eiendom, som er dets
intellectuelle Fædrenearv og et af de meest charakteristiske Udtryk af hele
Folkets Tænkeevne. Kategorierne komme vel først til Individet som en Tradition i
Sproget, men optages derfor ikke af det som noget Fremmed. Naar denne Tradition
tilegnes af et velbegavet Individ, finder det tvertimod deri en Form for sin
Selvtænkning, som det bruger saa sikkert, som om det selv havde produceret den.
Det tænker selvstændigt i de Former, som høre Stammen til, som ikke meer er det
ene Individs Eiendom end det andets, fordi det er en Form for den almindelige
Fornuft, der kun har Virkelighed, idet den gjentager sig i de særlige Individers
Tænkning. \opage{198}Individet tænker selv, naar det tænker i sin Stammes
Tankeformer; thi dets Fornuftliv bestaaer kun i Sammenhæng med et Selskabs
Fornuftliv, hvis store Strøm ogsaa gaaer igjennem dets individuelle Bevidsthed.

I Anvendelsen af Kategorierne og i Reflexion over de Kategorier i dets
Bevidsthed, der ikke have tabt deres umiddelbare Betydning, føler det tænkende
Subject sig i sit eget Element, meer end i Kundskaber om Tingenes sandselige
Bestemmelser: derom maa en umiddelbar Bevidsthed kunne overbevise hvert tænkende
Subject. Naar det modtager den Kundskab, at en Ting er rød eller suur, da er
denne Kundskab udvortes for det i en anden Betydning, end naar det erkjender, at
eet Phænomen er Aarsag og et andet Virkning. Tænkningen føler sig anderledes ved
sig selv, naar den reflecterer over Begreberne Eenhed, Fleerhed, Grund, Følge
o.\ s.\ v., end naar den betragter de sandselige Prædicater, som tillægges de
tilværende Ting. De først nævnte Begreber ere nemlig Tænkningens egne Producter;
den gjenkjender i dem de faste Spor af sin egen Virksomhed; den føler sig her paa
en anden Maade hjemme, end naar den har at gjøre med Forestillinger og Begreber,
som fremstille blot sandselige Affectioners Indhold.

Kategorierne kunne saavel fremtræde paa en vulgær Maade, som med en høiere
Betydning. For det Første kunne de ligge til Grund for Tænkningen i den blot
naturlige Bevidstheds Periode. De bruges da, eftersom Raden kommer til dem, hver
for sig, efter en sund Tact af det tænkende Subject, medens det paa en sandselig
og forstandig Maade opfatter sin Erfaringskreds eller i en populær Form udtaler
sin Religionskundskab. Det formaaer ei at danne sig Forestillinger om de
Gjenstande, der i Omverdenen fremstille sig for dets Sandser, uden at dets
Anskuelse bestandig træder i Forening med en Tænkning, der gaaer ud fra
Kategorierne som sine faste Forudsætninger. Her \opage{199}vise de aprioriske
Begreber sig netop som aldeles klare og forstaaelige for Tænkningen; thi skjøndt
Begreberne Tilværelse, Aarsag, Virkning o.\ s.\ v. her idelig komme til
Anvendelse, føles paa dette Standpunct ingen Trang til disses Forklaring: der
spørges i denne Periode aldrig om, hvad Tilværelse, Aarsag, Virkning betyder;
det forstaaer sig her af sig selv; thi man har i disse Begreber med sin
Forstaaelses Midler at gjøre. Man tænker blot paa at forstaae Tingene, ei at
forstaae sin egen Forstaaelse. Saaledes har man tidt her i Danmark, ja selv for
gandske nylig, ført en ivrig Strid om den rette Grændse for Udenadslæren og
Forstaaen i den populære Religionsunderviisning, uden at de stridende Parter
have yttret mindste Tvivl om, at de forstode hinandens Begreb om Forstaaelse.
Tvertimod have de med Sikkerhed lagt dette Begreb til Grund som et fast, fælleds
Udgangspunct for deres Undersøgelser. Selv i den Periode, da Kategoriernes System
forstyrres ved Reflexioner og tildeels fastholdes ved isolerede Definitioner,
kunne de dog derfor staae i den vulgære Tænknings Tjeneste. Dette gjøre de
nemlig, saalænge de som isolerede, faste Udgangspuncter bruges under Opfattelsen
af den givne Virkeligheds Phænomener.

En dybere Betydning faaer en Kategorie, naar den fattes som en Bestemmelse for al
Virkelighed. En saadan universel Betydning finder man, at snart een, snart en
anden Kategorie har vundet i Videnskabens Historie. Undertiden har nemlig et
eller andet Individ pludselig indseet et saadant Begrebs Gyldighed for alt
Værende, naar det tænkes i sin Totalitet. Det har ved en saadan eenlig
Bestemmelse omfattet den hele Virkelighed i sin Tænkning og holdt for, at det
derved har tilfulde begrebet den i sin Sandhed. Saaledes have flere Philosopher
med Enthusiasme forsvaret den Sætning, at Alt var Eet, og have levet i den
Overbeviisning, at de i dette Begreb om Eenhed havde fundet Nøglen til hele
Tilværelsens Gaade. Men en saadan \opage{200}eensidig Lære har gjerne fremkaldt
en ligesaa kraftig Fremstilling af den for en udvortes Betragtning aldeles
modsatte Sætning, at det i Sandhed Værende er en uendelig Mangfoldighed. (Vi
vælge nemlig her et Exempel, der, som bekjendt, allerede blev fremført af den
xenophontiske Sokrates, da han med faa Ord fremhævede den store Modsigelse, der
fandt Sted mellem Philosophernes Principer). --- Men, naar en Tænker holder for,
at han i en eenlig Kategorie har opdaget Virkelighedens sande Begreb, da ligger
der i hans Mening baade Sandhed og Vildfarelse. Sandhed er det, at en eenlig
Kategorie saaledes kan tænkes, at man deri har et Moment af Virkelighedens
Begreb, som kan lægges til Grund for en Verdensanskuelse. Urigtigt er det
derimod, at han frakjender andre Kategorier den samme universelle Betydning. Den
udtrykkelig udtalte Kategorie betegner kun een af den aprioriske Erkjendelses
Sider; men den, som i et saadant eenligt Ord finder en tilfredsstillende Form
for sin Erkjendelse, mener ved sit Ord meget Meer, end derved er udtalt. Hvis
han var i Stand til at bringe sin Erkjendelses Eensidighed til klar Bevidsthed,
vilde han med det Samme være kommet til den Indsigt, at han maatte skride frem
til en modsat Bestemmelse, der havde ligesaa universel Betydning. Men de
hinanden modsatte Bestemmelser forenes i en tredie, rigtigere Kategorie, hvori
de to foregaaende begge bevares i deres relative Gyldighed. Den tredie, fyldigere
Kategorie bliver ved fortsat Tænkning ligeledes erkjendt for eensidig, gaaer
over i sin Modsætning og danner derpaa i Forening med den en ny, fuldstændigere
Kategorie.

(Ved disse Par Ord er en foreløbig Forestilling givet om Ontologiens Methode,
der vil vinde Klarhed ved Fremstillingen af Videnskaben selv. Desuden agter jeg,
deels under Foredragets Løb, deels ved Slutningen deraf, at tillade mig endeel
oplysende Reflexioner, i den Hensigt, at bringe Methoden til klar Bevidsthed).

Ved en saadan Tankebevægelse kommer man til et System \opage{201}af aprioriske
Begreber, hvilket, naar Begyndelsen er tænkt som Begrebet om al Virkelighed, i
en meer udviklet Form beholder den universelle Betydning, eller fremstiller en
Erkjendelse „\textit{sub specie æternitatis}.“ I et saadant System af Begreber
udtales en oprindelig Erkjendelse af Virkeligheden, fordi deri bestemmes, hvad
der altid og nødvendig maa tænkes ved Virkeligheden, ihvordan den end forresten
fremstiller sig. En saadan apriorisk Tankekreds, som udgjør Ontologiens Indhold,
giver det Resultat, at Personlighed er Virkeligheden tænkt i sin Sandhed, eller
at Intet er med Sandhed at kalde virkeligt, uden hvad der har Personlighed. Men
kun i Bevægelsen gjennem den hele Cyclus af eenlige Bestemmelser bliver det
tydeliggjort, hvad Virkelighed er og hvad Personlighed er; thi paa dette Sted
have vi kun forudsat, at ethvert af disse Ord selv medfører en umiddelbar
Betydning. Ved en saadan apriorisk Tankebevægelse opnaaes en dobbelt Fordeel.
For det Første bliver Virkelighedens aprioriske Begreb fuldstændig forklaret.
Der gjøres nemlig Rede for, hvorledes Virkeligheden maa tænkes altid eller under
Evighedens Form: der gjøres saaledes Rede derfor, at det først ved Slutningen af
Foredraget er udtalt, hvad man forstaaer ved den Sætning, at Virkelighed er
Personlighed. Denne Sætning maa ved Foredragets Slutning have vundet en anden
Betydning end den, hvormed den her uden Forberedelse kan udtales. --- For det
Andet blive de eenlige Kategorier med det Samme forklarede, naar de under
Foredragets Løb defineres paa deres Plads og ved deres Plads i den hele
Sammenhæng. For den, der, efterat have gjennemtænkt hvert Grundbegreb efter dets
Betydning i den ontologiske Tankekreds, atter vender sig til den empiriske og den
praktiske Tænkning, have disse Grundbegreber, skjøndt de bruges i de samme
Forbindelser, hvori den naturlige Sprogtact valgte dem, faaet en dybere
Betydning; thi selv til det daglige Livs Kreds \opage{202}medbringer han nu det
Lys, som Ontologien har gjennemtrængt dem med.\footnote{Ontologien er Læren om
Tænkningens og Tilværelsens evige Form: den fremstiller nemlig de almindelige
Bestemmelser, som gjelde for Alt, hvad der er, hvilket Alt tillige, som
Videnskabens følgende Fremstilling skal vise, medfører de Bestemmelser, hvormed
vi hist og her i det Foregaaende ved Anticipation have betegnet det, nemlig
Tilværelse, Virkelighed. Ved Udviklingen og Forklaringen af de bestandig
gjeldende Bestemmelser for det Virkelige bliver tillige Gud og Verden og Guds
Forhold til Verden tænkt i abstracte Begreber ɔ: tænkt i Almindelighed, uden at
derfor hele Virkelighedens concrete Indhold optages i den ontologiske Tænknings
Form. Ligesom Tænkeren i en enkelt Kategorie tidt har tænkt Alt under Evighedens
Form, men saaledes at hans Tænkning paa en abstract Maade deri omfatter al
Virkelighed, saaledes er ogsaa hele Tilværelsen abstract tænkt i det hele System
af aprioriske Begreber, som vel med objectiv Gyldighed bestemmer al Virkelighed,
men ei fremstiller hele Tilværelsens Rigdom. (Uden at indrømme noget Individ
vilkaarlig Magt til at bestemme den Talebrug, som skal bruges i Philosophiens
frie Republik, bruge vi dette mangetydige Ord --- abstract --- her i en
Bemærkelse, som har tilstrækkelig Hjemmel i Videnskabens Historie. Det bruges i
Videnskaben snart om det Residuum af Anskuelsen, der beholdes i Forestillingen,
snart, som her, om de ontologiske Begreber i deres Totalitet, snart om en eenlig
Kategorie, hvori færre Tankebestemmelser indeholdes, forsaavidt den
sammenlignes med en indholdsrigere Kategorie. Alt Aristoteles betegner ved
τὸ χωριστόν saavel en af de for sig tænkte Kategorier, som det af
særlige Sandsegjenstande afdragne Begreb). --- [Herved er i Randen af
Forfatterens Haandskrift bemærket: „Denne § maa nærmere overveies, især
Talebrugen: Ontologie og Metaphysik.“]}

Ontologien udtaler, som sagt, et Hovedmoment af den \opage{203}videnskabelige
Erkjendelse: men dennes andet Hovedmoment er Erfaringen. Kun den Erfaring, der
har det gjennemarbeidede System af Kategorier til sin Forudsætning, eller den
aprioriske Tænkning, der er forenet med Sympathie og virkelige Personers
Erkjendelse, er menneskelig Erkjendelses Energie. Den eensidige aprioriske
Tænkning af Tænkningen er ei guddommelig, men i Værdi under den sande
menneskelige. Erfaring uden den rene Tænkning er blind, og apriorisk Begreb uden
personlig oplevet Erfaring er blot formel Begribning eller
Schematismus.\footnote{Dette Stykke kan maaskee meer udvikles, især Forholdet
imellem reen Tænkning og Erfaring. \emph{Forf.s Marginalnote.}}

Ontologien har dette dobbelte Formaal, deels at bringe de abstractere
Bestemmelser, som udgjøre de concretere Kategoriers Indhold, til Bevidsthed,
deels at vise, hvorledes hine Bestemmelser med Nødvendighed forene sig i de
concretere, saa at Tankebevægelsen fortjener Navn saavel af analytisk som
synthetisk; thi netop ved den adskillende Tænkning, Adskillelsen af en eenlig
abstract Bestemmelse, bliver dens Uadskillelighed fra den modsatte Bestemmelse,
hvorved den integreres, allertydeligst indseet. Men her opstaaer det Spørgsmaal,
hvorvidt Analysen skal gaae; thi En og Anden synes at antage det for en afgjort
Sag, at den fuldkomneste Tænkning her adskiller de fleste Momenter ---
naturligviis under den Forudsætning, at Synthesen bestandig følger med. At
Adskillelsen her ikke støder paa nogen absolut Grændse, seer man let, da den
finere Tænkning bestandig har det i sin Magt, at indskyde flere Led i de discrete
Begrebers Række. Men, da Tænkningen dog kun er ægte (sand) Tænkning ved bestemt
at udtales, kan den af Mangel paa passende Ord let komme til at slæbe sig frem
med et Skinliv. Den Tænkende kan let komme til at tabe sig i Haarkløverier, hvor
Ordet gandske træder i den \opage{204}formeente Tankes Sted. Ligesom Talen ved
at efterligne Tankebevægelsens Form kan fortsættes, efterat Tænkningen har
ophørt at bevæge sig deri, kan der ogsaa paa Overgangen til den tomme Talen
gives en mat og halvdød Talen. Sit egentlige, friske Liv har saavel den
philosophiske Tænkning som anden Tænkning kun, medens Ordets umiddelbare
Betydning i Modersmaalet ei reent er forladt. Deraf har den græske Philosophie
sine friske Farver; thi den er voxet frem af Folkets (Stammens) populære
Tænkning, og den bæres af Ordenes umiddelbare Betydning som af en kraftig Rod.
Svagere bliver den Tænkning, som gaaer for sig i fremmede Folkestammers Ordformer
og træder frem i en broget Dragt af udenlandske Talemaader. Dog kunne de endnu
for den Lærde, som har sat sig ind i fremmede Folkeaanders Talebrug, have meget
deriveret Liv tilbage. Men, naar individuel Vilkaarlighed i Ordenes Dannelse
faaer altfor stort Raaderum, taber Talen alt meer og meer af sin Betydning, fordi
den udbreder sig i Luften, løsreven fra sin Rod og sit Livsprincip. Længst
vedbliver den som en blot subjectiv Form at have Gyldighed for Tænkeren selv, men
tilsidst bliver den eensomme Tænkning i selvgjorte Former til et blot Ordmageri.
Heri ligger en af Grundene til, at den analyserende Tænkning kan gaae for vidt;
den kan nemlig nærme sig for meget til det Punct, da Ordet bliver deels et
vilkaarligt, blot subjectivt Udtryk for Tanken, deels gandske tankeløst. Den
gamle Sætning, \textit{qui bene distinguit, bene docet}, er derfor temmelig
eensidig, med mindre man lægger særdeles Eftertryk paa Tillægget \textit{bene}.
Ved dette \textit{bene} maa da deels udtrykkes, at Distinctionen maa bestandig
gaae Haand i Haand med Conjunctionen, deels, at Analysen ei maa fortsættes, til
den bliver en indholdsløs, blot skrømtet Tænkning.

Men, at en Fremstilling af Philosophien (Ontologie) har overskridt de rette
Grændser for Begrebets Articulation og for\opage{205}vildet sig i
Spidsfindigheder, kan kun umiddelbart dømmes af den skjønsomme Kjender; der lader
sig naturligviis ikke føre noget videnskabeligt Beviis derfor. Den rette Grændse
bliver her efter Individualiteternes (Folkets) og Tidsaldrenes Fordringer af
høist foranderlig Natur, saa at heri ligger en af de mange Grunde, som nøder
Philosophien til at træde frem i en Mangfoldighed af individuelle Former. Den
Fremstilling af Philosophien, hvorved et Individ uden Affectation finder sig
tilfredsstillet, og hvori det bevæger sig frit, er snart meer concentreret, snart
meer udviklet; men een eneste for Alle gyldig Norm kan aldrig findes. Her
opstille vi kun den almindelige Sætning, at Sondringen af de abstracte
Bestemmelser, som i et Begreb ere forenede, kan drives for vidt. Afhandlinger,
hvori denne Feil begaaes, tildrage sig derfor liden Opmærksomhed i Literaturen:
de have nemlig kun Interesse for Skribenten selv, men ikke for Læseren. En
saadan overdreven Analyse foretages ikke sjelden af Hegelianere, som efter Hegels
egen Methode fortsætte Sønderlemmelsen af de Begreber, som han bliver staaende
ved. `Hvor livløs er ikke Fremstillingen i flere Værker af Hinrichs! Selv Hegel,
som dog havde meest Anledning til at glæde sig ved den, beder i sine Breve
Forfatteren at lægge større Vind paa Klarhed, og tilstaaer, at det ofte falder
ham selv temmelig besværligt at følge hans Foredrag.

Ligesom Philosophiens Fremstilling tidt taber sin Intensitet derved, at Begrebet
bliver altfor meget opløst i sine abstracte Bestanddele og saa at sige optrævlet
i tynde Traade, lider den ogsaa tidt af den modsatte Mangel og taber derved sin
objective Gyldighed; vi tænke herved paa det overfladiske Foredrag, som gjør
altfor lange Skridt i sin Dialektik, overspringer for mange Mellemled og derved
bliver uforstaaeligt for Læseren. Men det rette Maal beroer ogsaa her paa Tiden
og de bestemte Forhold, hvorunder Fremstillingen er bleven til. --- Saaledes er
det bestemte \opage{206}Forhold mellem Intension og Extension i den philosophiske
Tænknings Formation en af de Maader, hvorpaa Nationalaanden udtaler sig hos et
Folk. --- ---

\bigskip
\begin{center}
{\large\scshape Første Capitel.}\\[0.4em]
{\large\bfseries Læren om de enkelte Begreber.}
\end{center}
\begin{center}\rule{0.25\textwidth}{0.4pt}\end{center}
\begin{center}
{\bfseries Første Afsnit.}\\[0.3em]
{\bfseries Begyndelsesbegreberne.}
\end{center}
\medskip

Ved Reflexioner over Erkjendelsens Natur kan den naturlige Vished om Tænkningens
objective Realitet i den Grad forstyrres, at al Sandhedserkjendelse synes at
forvandle sig til et betydningsløst Spil med blot subjective Begreber. Derved
opstaae saadanne Skikkelser af Bevidstheden som Idealisme, Skepticisme og
Nihilisme, der have deres fælleds Grund deri, at Tænkeren først ved en voldsom
Abstraction aldeles adskiller Tænkning fra Tilværelse, og dernæst forgjeves seer
sig om efter Grunde for deres Overeensstemmelse.

Men ved at fortsætte den Tankegang, der betragter Tænkningens Grundbegreber som
blot subjective Former, mærker man, at den umulig lader sig gjennemføre med
Consequents. Den lader sig aldeles ikke udtale uden under Forudsætning af visse
Grundbegrebers reale Gyldighed, som man kan betjene sig af som Midler til dens
Fremstilling. Men, naar man har gjort sig dette klart, er den blot subjective
Tænkning overvunden; den Selvmodsigelse, som den tilsidst indvikler sig i, er
bleven bragt til Bevidsthed, og Veien er banet for et Fremskridt i Erkjendelsen.

De Begreber, hvorved Tilværelse tænkes, kunne nemlig \opage{207}umulig alle
tilsammen tabe deres Betydning for Tvivleren. Det erfares desto tydeligere, jo
mere man nærmer sig med sin skeptiske Reflexion til de meest abstracte blandt
dem. Vi ville strax henvende Opmærksomheden paa det meest indholdsløse
Grundbegreb, som forefindes i vor Bevidsthed, nemlig Begrebet Væren eller at
være; thi deri møde vi netop en uoverstigelig Grændse for idealistiske Tvivl og
Grublerier. Den Mening, at Begrebet Væren ikke har nogen objectiv Gyldighed, har
Ingen consequent forsvaret, og den kan Ingen tydelig udtale; thi hvis man ikke
forudsætter dette Begrebs Realitet, kan man om Intet paastaae, at det er, eller
at det ikke er. Deraf kommer den tidt udhævede Modsigelse, hvori Kants
Fornuftkritik er indviklet. Efter Systemets Plan maatte Væren ansees for et
subjectivt Begreb, der blot gjeldte for Phænomenernes Verden; men desuagtet
tales der idelig om det, der i sig selv er, uden at der gjøres mindste Rede for,
hvorfra Begrebet herom hentes, og hvorfra det faaer sin Realitet. Men det er
Sagens Natur, som her har gjort sig gjeldende mod en kunstlet Tankegang: man kan
Intet paastaae og Intet benegte uden at tillægge Begrebet Væren meer end
subjectiv Betydning; ja uden at antage dets Realitet har man ikke Midler til at
frakjende det sin Realitet. --- Dette Begreb, som medfører saadan uovervindelig
Vished om Realitet, er høist skikket til at gjøre Tænkningens uopløselige
Sammenhæng med Væren ret evident, og dette Begreb have vi, som sagt, valgt til
Udgangspunct for Ontologien.

\emph{A}~~\emph{Væren.} Da Ontologien her foredrages uden sin videnskabelige
Indledning, tages Begrebet Væren her som umiddelbart. Der forudsættes, at Ordet
medfører den Betydning, hvori det her skal tages, men ved Fremgangen til de
concretere Begreber, til hvilke dette Begyndelsesbegreb dialektisk vil udvikle
sig, vil \opage{208}denne Betydning træde bestemtere frem, idet Overgangen selv
her, ligesom ved alle de følgende Kategorier, vil kaste Lys tilbage paa det
Foregaaende.

Tanken Væren er Begyndelsespunctet i den egentlig philosophiske eller productive
Tænkning, og det er Videnskabens Formaal, at lade et heelt apriorisk Tankesystem
fremstaae ved dette Puncts Udvidelse. Idet Begrebet Væren for Tænkningen udvikler
sig til indholdsrigere Begreber, kan Tænkningen med Sandhed siges selv at
frembringe sit Indhold, og det Indhold, som den selv frembringer, er en Cyclus af
Bestemmelser for al Virkelighed, hvilke have Almeengyldighed og Nødvendighed. ---
Undertiden bestrides den Paastand, at man ved en apriorisk eller, som man nu tidt
hører, immanent Tænkning kan udvikle Tænkningens og Tilværelsens almindeligste
Bestemmelser saaledes, at Tænkningen efter Abstraction fra alt Indhold med fri
Selvstændighed reproducerer sig selv i den sande Form, der tilkommer den.
Indsigelserne mod denne Lære beroe aldeles paa Misforstaaelse. Men denne
Misforstaaelse lader sig hæve ved et Par Ord, som tillige kunne bidrage til at
drage Noget af den fornemme Dunkelhed bort fra den immanente Tænknings Begreb.
Naturligviis indrømme vi, at Tænkningens Grundbegreber udvikle sig under
Erfaringen, og at de vigtigste Kategorier, som udgjøre Ontologiens Indhold,
førend Ontologiens Idee kan fattes, maae forefindes i Tænkernes Bevidsthed. Den
omtalte Production er derfor i en vis Betydning en grundigere Reproduction,
hvorved man deels tydeliggjør sig det Forhold, der finder Sted mellem hine
Grundbegreber, af hvilke man vel før havde gjort en umiddelbar Brug, deels, hvad
der er det Vigtigste, idet man ordner dem, finder dem at være sammenhørende
Brudstykker af en oprindelig menneskelig Grunderkjendelse. Ved denne forandrede
Stilling og Betydning, som Kategorierne faae efter hiin Reproduction, berettiges
man altsaa virkelig til at betragte den som en \opage{209}Production af noget
Nyt. --- Desuden er Ontologen ingenlunde indskrænket til den Cyclus af
Kategorier, der overleveres ham som udtrykkelig benævnte i Sproget. Medens han
paa en methodisk Maade befinder sig over Kategoriernes Sammenhæng, kan det vel
være muligt, at Videnskaben kræver en anden Articulation, end den, som det
sociale Livs Trang har foranlediget. I enhver grundig Ontologie findes derfor
Kategorier, som ei kjendes udenfor det videnskabelige Sprog. Men her er det, at de
luftige Formationer af Ord og Formler saa let kunne skaffe sig Indpas, idet
Herskesygen og den subjective Vilkaarlighed gjør sig gjeldende under Valget af
Kunstord. Først, naar denne Mislighed tager Ende, vil Philosophiens Studium
blive ret glædeligt, fordi den menneskelige Tænkning kan komme til at bevæge sig
i Former, der stemme med Naturens og Skjønhedens Love. Da først ville
Philosophiens sande Præster komme til at føre dens Ord. --- Hvo tvivler vel om,
at nutildags de, som egentlig af Naturen ere kaldede til philosophisk Forskning,
finde sig udelukkede derfra ved de gjeldende Philosophers Nærlighed, Mangel paa
Fremstillingsgave og fremherskende Laugsaand. Philosophisk Genialitet er vist
ikke det fornemste Requisit for den, der skal kunne orientere sig i den nu
gjeldende tydske Philosophies factiske Tilstand.

Som Begyndelse paa al Tankebestemmelse\footnote{Her har Forfatteren antegnet
følgende: „Tildeels skrevet under min Sygdom i dette Aar (1838) og begyndt under
asthmatiske Anfald i Fjor, efterat et andet mislykket Udkast var casseret.“} maa
Begrebet Væren af alle være det meest abstracte; der lade sig ikke adskille
forskjellige Bestemmelser deri som dets Indhold; det er aldeles enkelt og
uopløseligt. Derfor er det umuligt at give nogen Definition derpaa i den
sædvanlige Betydning af Ordet Definition. Dertil udfordredes nemlig, at der skulde
gives et meer abstract \opage{210}Begreb, hvorunder det kunde subsumeres; men da
maatte dette Begreb betragtes som et rigtigere Udgangspunct for den ontologiske
Tænkning end Begrebet Væren. Begrebet Væren bliver kun defineret ved sin Plads i
Kategoriernes hele System. En Følge af Værens abstracte Natur er dets
Almeengyldighed. Alt, hvad der lader sig bestemme ved nogen anden Kategorie, som
Aarsag eller Følge, eller Substants o.\ s.\ v., lader sig ogsaa bestemme ved
dette alleralmindeligste Begreb, som ligger til Grund for dem alle. Naar en af de
concretere Kategorier prædiceres om et Subject, er dermed Væren indirecte
prædiceret derom; derimod er med den blotte Væren ei en eneste af de øvrige
Tankebestemmelser udtalt om et Subject.

Da Væren paa en abstract Maade omslutter al Virkelighed og indeholdes i alle
andre Tankebestemmelser for Virkeligheden, bliver Begrebet tænkt urigtigt, hvis
det opfattes som en Modsætning til Tænkning. Saaledes kan nemlig Ordet, efter sin
umiddelbare Betydning, ogsaa forstaaes. Men efter den altomfattende Betydning,
hvori Væren her skal tages, gjelder det saavel om alt Subjectivt som om alt
Objectivt. Ogsaa Tænkningen falder under Bestemmelsen Væren.

Som Overgang til reel Tænkning eller Begyndelse dertil er Væren paa eengang baade
Bestemmelse og Ikke-Bestemmelse. Denne Modsigelse gjør, at man ikke kan blive
staaende ved dette Begreb, men drives derfra til andre Begreber, hvorved det
fuldstændiggjøres. Den Trang til nærmere Bestemmelse, som følger med Begrebet
Væren, ville vi endnu ved et Par forskjellige Vendinger stræbe at bringe til
klarere Bevidsthed.

Væren kan betragtes som det blotte Copula med et aldeles ubestemt Subject, der
søger sit Prædicat. Det udsiger hverken, at Noget er, eller hvad det er.
\textit{x} er \textit{x} er den Formel, hvori denne fuldkomne Mangel paa
Bestemthed kan udtales. Ved denne \opage{211}Reflexion, at vi her af Sætningens
Form kun have det blotte Copula --- er ---, bliver det ret anskueligt, at vi have
at gjøre med en ufuldført Tanke, der kræver sit Supplement.

Den aldeles ubestemte Væren er endnu Intet. Saaledes er man selv efter den
umiddelbare Sprogbrug berettiget til at udtrykke sig. Hvad der ikke kan bestemmes
ved noget Prædicat, hvad der ikke er Noget, det er Intet. (Dette er een af de
flere Bemærkelser, hvori Ordet Intet kan tages). Ved den Betragtning, at Væren er
Intet, eller rettere, at den endnu for os er Intet, er den ontologiske
Tankebevægelse egentlig først begyndt. Den forudsætter nemlig ligesom enhver
Bevægelse to Puncter: et, hvorfra Bevægelsen gaaer ud, og et, til hvilket den
gaaer hen. Ved denne Bevægelse bliver det her bragt til bestemt Bevidsthed, at
der i Begrebet Væren ligger et Savn eller en Mangel. Intet betegner her den
ledige Plads for et Prædicat. Den rene Væren er Intet, siger det Samme som: at
være er endnu Intet, med mindre derved tænkes at være Noget. \emph{Noget} er
saaledes den nye Kategorie, som træder i Stedet for Intet, og først ved at udtale
den Sætning: „\emph{Væren er Noget}“, har jeg gjort et klart Fremskridt i
Tænkningen.

Det er blevet Skik og Brug efter Hegels Autoritet at betragte den rene Væren som
den Kategorie, der ligger til Grund for Eleaternes Verdensanskuelse. Uden at
afgjøre det Spørgsmaal, om hiin Philosophie med Grund jævnføres med dette Sted af
Ontologien, (hvilket maaskee fra et strengt historisk Standpunct maatte
benegtes), bemærkes i Forbigaaende at τὸ ὄν hos Eleaterne --- i Stedet for
τὸ εἶναι --- i det Mindste ikke nøiagtig svarer til den rene Væren eller det at
være, da den Hypostasering, som ligger i det græske Ord, gjør Begrebet meer
concret, end det her efter sin Bestemmelse skulde være. Det danske Væren er her
meer adæqvat end de tilsvarende Ord hos flere \opage{212}andre dannede Nationer.
Det franske \textit{être} er f.\ Ex.\ mere tvetydigt saavelsom det tydske
\emph{Seyn}, der, i det Mindste med den bestemte Artikel, ogsaa bruges for at
betegne Indbegrebet af Alt det, som er.

\emph{B.}~~\emph{Noget.} Det rette Fremskridt fra den ubestemte Væren er det, som
gjøres til Bestemthed i Almindelighed. At være er at være bestemt paa en eller
anden Maade, eller Alt, hvad der er, er bestemt, er Noget --- ---

\begin{center}\rule{0.32\textwidth}{0.4pt}\end{center}

\medskip
\begin{center}\footnotesize\textit{[Brudstykket afbrydes her.]}\end{center}

\end{document}
